\chapter{Inequalities between Matrices}\label{linear:ch:inequalities}

\begin{center}\itshape
All analysts spend half their time hunting through the literature for inequalities which they want to use and cannot prove.
\upshape(G.~H.~Hardy, \emph{Prolegomena to a chapter on inequalities}, 1929)
\end{center}

In 1893 Jacques Hadamard, then a young teacher at the Lyc\'ee Buffon in Paris, asked how large the determinant of an $n\times n$
matrix can be when its entries are bounded by $1$ in absolute value. The determinant is the volume of the parallelepiped spanned by
the columns, each column has length at most $\sqrt n$, and a box is largest when its edges are perpendicular: so $|\det A|\le n^{n/2}$,
with equality exactly for matrices of $\pm1$ whose columns are orthogonal. Hadamard proved this, built such matrices for $n=12$
and $n=20$, and observed that they can only exist when $n$ is $1$, $2$ or a multiple of $4$. Whether they exist for \emph{every}
multiple of $4$ is the \emph{Hadamard conjecture}. It is still open after more than a century; the smallest unsettled order is
$668$. (Volume~14, Appendix~G, tells that story.)

Hadamard's inequality is the ancestor of a whole family of inequalities between matrices, and of a whole way of thinking: matrices
can be \emph{compared}, $A\succeq B$, much as numbers can, as long as one is careful about what survives the passage from numbers to
matrices. Some familiar facts survive ($\sqrt{\phantom{t}}$ is increasing, $\log\det$ is concave, products of positive things are
positive, if the product is the right one), and some do not ($t\mapsto t^2$ is not increasing on matrices). This chapter proves
the most useful of these inequalities, with an eye on what they mean: volumes, ellipsoids, covariance matrices and kernels. It ends
with Gershgorin's discs, the cheapest eigenvalue bound there is, and with Hilbert's inequality, whose constant $\pi$ comes out of a
convex function and an integral.

\begin{think}
If $A$ and $B$ are symmetric $2\times2$ matrices with $A-B$ positive semidefinite, must $A^2-B^2$ be positive semidefinite?
Before computing, think of $A$ and $B$ as ellipses. What would a counterexample look like?
\end{think}

\section{The Loewner order}

Recall (Chapter~\ref{linear:ch:spectral}) that a Hermitian matrix $A$ is \emph{positive semidefinite}, $A\succeq0$, if $x^*Ax\ge0$ for all $x$,
equivalently if its eigenvalues are $\ge0$; and \emph{positive definite}, $A\succ0$, if $x^*Ax>0$ for $x\ne0$. For Hermitian $A,B$ we
write $A\succeq B$ when $A-B\succeq0$. This \emph{Loewner order}, named after Charles Loewner (Karel L\"owner), who studied it in 1934,
is a partial order, and most of its basic properties are one line each.

\begin{proposition}\label{linear:prop:loewner}
Let $A,B$ be Hermitian of size $n$.
\begin{enumerate}[label=(\roman*)]
\item If $A\succeq B$, then $X^*AX\succeq X^*BX$ for every $n\times m$ matrix $X$ (\emph{congruence preserves the order}).
\item If $A\succeq B$, then $\lambda_k(A)\ge\lambda_k(B)$ for every $k$; in particular $\operatorname{tr}A\ge\operatorname{tr}B$, and
$\det A\ge\det B$ if $B\succeq0$.
\item If $A\succeq B\succ0$, then $B^{-1}\succeq A^{-1}\succ0$.
\item $A\succeq B$ if and only if the ellipsoid $E_A=\{x:x^*Ax\le1\}$ is contained in $E_B$ (for $A,B\succ0$).
\end{enumerate}
\end{proposition}
\begin{proof}
(i) $y^*X^*(A-B)Xy=(Xy)^*(A-B)(Xy)\ge0$. (ii) By Courant--Fischer (Theorem~\ref{linear:thm:courantfischer}), $\lambda_k$ is a max-min of
Rayleigh quotients, and $x^*Ax\ge x^*Bx$ for every $x$. The determinant is the product of the eigenvalues, which are $\ge0$ when
$B\succeq0$. (iii) Let $C=B^{-1/2}AB^{-1/2}$, where $B^{-1/2}$ is the inverse of the positive square root (Exercise~\ref{linear:ex:sqrtpsd}).
By (i), $C\succeq B^{-1/2}BB^{-1/2}=I$, so the eigenvalues of $C$ are $\ge1$ and those of $C^{-1}=B^{1/2}A^{-1}B^{1/2}$ are $\le1$: $C^{-1}
\preceq I$. Applying (i) with $X=B^{-1/2}$ gives $A^{-1}\preceq B^{-1}$. (iv) If $A\succeq B$ and $x^*Ax\le1$, then $x^*Bx\le x^*Ax\le1$.
Conversely, if $x^*Bx>x^*Ax$ for some $x$, scale $x$ so that $x^*Ax<1<x^*Bx$ (if $x^*Ax=0$ then $x=0$, as $A\succ0$): this point is
in $E_A$ but not in $E_B$.
\end{proof}

So the Loewner order is the order of ellipsoids: bigger matrices, smaller ellipsoids. Figure~\ref{linear:fig:loewner} shows why squaring is
not monotone. With
\[
A=\begin{pmatrix}2&1\\1&1\end{pmatrix},\qquad B=\begin{pmatrix}0.9&0\\0&0.05\end{pmatrix},
\]
$A-B$ has positive trace and determinant $1.1\cdot0.95-1>0$, so $A\succeq B$; but $A^2-B^2=\bigl(\begin{smallmatrix}4.19&3\\3&1.9975
\end{smallmatrix}\bigr)$ has negative determinant. Squaring stretches $E_A$ along its long axis, which is tilted, and it pokes out of
the strip $E_{B^2}$.

\begin{figure}[ht]
\centering
\fig{linear-ch11-loewner}
\caption{Left: $A\succeq B$ means that the ellipse $E_A$ (blue) lies inside $E_B$ (red, here a long ellipse, clipped). Right:
$E_{A^2}$ is thinner and longer along a tilted axis, and leaves $E_{B^2}$: $A^2\not\succeq B^2$. Squaring is not monotone for matrices.}
\label{linear:fig:loewner}
\end{figure}

Square roots, however, are monotone.

\begin{theorem}[L\"owner--Heinz for the square root]\label{linear:thm:sqrtmonotone}
If $A\succeq B\succeq0$, then $\sqrt A\succeq\sqrt B$.
\end{theorem}
\begin{proof}
Let $S=\sqrt A$, $T=\sqrt B$, and let $\lambda$ be an eigenvalue of the Hermitian matrix $S-T$, with unit eigenvector $x$. Since
$(S-T)x=\lambda x$ and, taking adjoints, $x^*(S-T)=\lambda x^*$,
\[
0\le x^*(A-B)x=x^*\bigl(S(S-T)+(S-T)T\bigr)x=\lambda\,x^*Sx+\lambda\,x^*Tx.
\]
If $\lambda<0$, this forces $x^*Sx+x^*Tx\le0$; both terms are $\ge0$, so both vanish, and for positive semidefinite matrices
$x^*Sx=0$ implies $Sx=0$ (write $x^*Sx=\|\sqrt Sx\|^2$). Then $Sx=Tx=0$ and $\lambda x=(S-T)x=0$, a contradiction. So every eigenvalue
of $S-T$ is $\ge0$.
\end{proof}

\begin{remark}
A function $f$ on $[0,\infty)$ is \emph{operator monotone} if $A\succeq B\succeq0$ implies $f(A)\succeq f(B)$ in every dimension. Heinz
proved in 1951 that $t^p$ is operator monotone for $0\le p\le1$, and Loewner's theorem of 1934 characterises all operator monotone
functions: they are exactly the restrictions of functions holomorphic on the upper half-plane, mapping it to itself (Pick functions).
So $\sqrt t$, $\log t$ and $\frac t{1+t}$ are operator monotone, $t^2$ and $e^t$ are not. A deep theorem from a two-line definition.
\end{remark}

\section{Schur complements}

Block matrices are where most matrix inequalities are born. Let
\[
M=\begin{pmatrix}A&B\\B^*&C\end{pmatrix}
\]
be Hermitian, with $A$ invertible. The \emph{Schur complement} of $A$ in $M$ is $M/A=C-B^*A^{-1}B$.

\begin{proposition}[Schur complements]\label{linear:prop:schurcomplement}
With $X=\bigl(\begin{smallmatrix}I&-A^{-1}B\\0&I\end{smallmatrix}\bigr)$,
\[
X^*MX=\begin{pmatrix}A&0\\0&C-B^*A^{-1}B\end{pmatrix}.
\]
Consequently $\det M=\det A\det(M/A)$, and if $A\succ0$, then $M\succeq0$ if and only if $M/A\succeq0$ (and $M\succ0$ if and only if
$M/A\succ0$).
\end{proposition}
\begin{proof}
Multiply out: $MX=\bigl(\begin{smallmatrix}A&0\\B^*&C-B^*A^{-1}B\end{smallmatrix}\bigr)$, and $X^*$ subtracts $B^*A^{-1}$ times the first block
row from the second, which kills $B^*$. Since $\det X=1$, the determinants agree. The matrix $X$ is invertible, so by
Proposition~\ref{linear:prop:loewner}(i) applied with $X$ and $X^{-1}$, $M\succeq0$ if and only if $X^*MX\succeq0$, that is, if and only if
$A\succeq0$ and $M/A\succeq0$.
\end{proof}

The same trick with rectangular blocks gives a famous identity.

\begin{proposition}[Sylvester's determinant identity]\label{linear:prop:detIAB}
For $A\in M_{m,n}(\C)$ and $B\in M_{n,m}(\C)$, $\det(I_m+AB)=\det(I_n+BA)$. In particular $\det(I+uv^*)=1+v^*u$.
\end{proposition}
\begin{proof}
Compute the determinant of $N=\bigl(\begin{smallmatrix}I_m&-A\\B&I_n\end{smallmatrix}\bigr)$ in two ways. Eliminating with the top-left block,
\[
\begin{pmatrix}I_m&0\\-B&I_n\end{pmatrix}N=\begin{pmatrix}I_m&-A\\0&I_n+BA\end{pmatrix},
\]
so $\det N=\det(I_n+BA)$. Eliminating with the bottom-right block, $N\bigl(\begin{smallmatrix}I_m&0\\-B&I_n\end{smallmatrix}\bigr)=
\bigl(\begin{smallmatrix}I_m+AB&-A\\0&I_n\end{smallmatrix}\bigr)$, so $\det N=\det(I_m+AB)$.
\end{proof}

The identity turns an $m\times m$ determinant into an $n\times n$ one, which is a large saving when $n\ll m$; it also says that $AB$ and
$BA$ have the same non-zero eigenvalues, with multiplicity (replace $I$ by $tI$). Statisticians use it daily in the form of the
\emph{matrix determinant lemma}, and with it the Sherman--Morrison formula for $(A+uv^*)^{-1}$ (Exercise~\ref{linear:ex:sherman}).

\section{Determinant inequalities}

\begin{theorem}[Hadamard's inequality]\label{linear:thm:hadamard}
If $A\succeq0$, then $\det A\le a_{11}a_{22}\cdots a_{nn}$. If $A\succ0$, equality holds only for diagonal $A$.
\end{theorem}
\begin{proof}
If some $a_{ii}=0$, then $e_i^*Ae_i=0$ forces $Ae_i=0$, so $\det A=0$ and the inequality holds. Otherwise let $D=\operatorname{diag}(a_{ii}^
{-1/2})$ and $C=DAD$, which is positive semidefinite with unit diagonal. Its eigenvalues $\mu_i\ge0$ have sum $\operatorname{tr}C=n$,
so by the inequality of arithmetic and geometric means
\[
\det C=\prod\mu_i\le\Bigl(\frac1n\sum\mu_i\Bigr)^n=1,
\]
which is $\det A\le\prod a_{ii}$. Equality in AM--GM forces all $\mu_i=1$, so $C=I$ (a Hermitian matrix with all eigenvalues $1$ is
the identity), and $A=D^{-2}$ is diagonal.
\end{proof}

\begin{corollary}[Hadamard's volume inequality]\label{linear:cor:hadamardvol}
For any $n\times n$ matrix $X$ with columns $x_1,\dots,x_n$, $|\det X|\le\|x_1\|\cdots\|x_n\|$, with equality if and only if the
columns are pairwise orthogonal or one of them is zero. In particular a matrix with entries in $[-1,1]$ has $|\det X|\le n^{n/2}$.
\end{corollary}
\begin{proof}
Apply Theorem~\ref{linear:thm:hadamard} to the Gram matrix $A=X^*X\succeq0$, whose diagonal entries are $\|x_i\|^2$ and whose determinant is
$|\det X|^2$. If no column is zero and $X$ is invertible, equality forces $X^*X$ to be diagonal, that is, orthogonal columns; if $X$ is
singular with non-zero columns, the inequality is strict.
\end{proof}

A matrix of $\pm1$ entries reaching $|\det|=n^{n/2}$ is a \emph{Hadamard matrix}: $HH^{\mathsf T}=nI$. The Walsh--Hadamard matrices of
Chapter~\ref{linear:ch:symmetry} are examples for $n=2^k$, and Exercise~\ref{linear:ex:hadamard4} shows why $4\mid n$ is necessary for $n\ge3$.

Hadamard's inequality has a block version, with the same proof idea organised by Schur complements.

\begin{theorem}[Fischer's inequality]\label{linear:thm:fischer}
If $M=\bigl(\begin{smallmatrix}A&B\\B^*&C\end{smallmatrix}\bigr)\succeq0$ with square diagonal blocks, then $\det M\le\det A\det C$.
\end{theorem}
\begin{proof}
If $A$ is singular, then $M$ is singular (a vector $x$ with $x^*Ax=0$ gives $(x,0)^*M(x,0)=0$, hence $M(x,0)=0$), and the inequality
reads $0\le0\cdot\det C$. If $A\succ0$, then $M/A=C-B^*A^{-1}B\succeq0$ by Proposition~\ref{linear:prop:schurcomplement}, and $C\succeq M/A$
because $B^*A^{-1}B\succeq0$. By Proposition~\ref{linear:prop:loewner}(ii), $\det C\ge\det(M/A)$, so $\det M=\det A\det(M/A)\le\det A\det C$.
\end{proof}

Iterating Fischer down to $1\times1$ blocks gives Hadamard again. Our next inequality is about sums.

\begin{theorem}[Minkowski's determinant inequality]\label{linear:thm:minkowskidet}
For $A,B\succeq0$ of size $n$,
\[
\det(A+B)^{1/n}\ge\det(A)^{1/n}+\det(B)^{1/n}.
\]
\end{theorem}
\begin{proof}
If $A$ and $B$ are both singular there is nothing to prove; by symmetry assume $A\succ0$. Both sides are multiplied by $|\det X|^{2/n}$
when $A,B$ are replaced by $X^*AX,X^*BX$, so with $X=A^{-1/2}U$ for a suitable unitary $U$ we may assume $A=I$ and $B=\operatorname{diag}
(d_1,\dots,d_n)$, $d_i\ge0$ (spectral theorem). We must show $\prod(1+d_i)^{1/n}\ge1+\prod d_i^{1/n}$. By AM--GM,
\[
\prod_i\Bigl(\frac1{1+d_i}\Bigr)^{1/n}+\prod_i\Bigl(\frac{d_i}{1+d_i}\Bigr)^{1/n}\le\frac1n\sum_i\frac1{1+d_i}+\frac1n\sum_i\frac{d_i}{1+d_i}=1,
\]
and multiplying by $\prod(1+d_i)^{1/n}$ gives the claim.
\end{proof}

So $A\mapsto\det(A)^{1/n}$ is superadditive, and, being homogeneous of degree $1$, concave on the cone of positive semidefinite matrices.
The logarithm of the determinant is concave too, and this is the form statisticians and information theorists use.

\begin{theorem}[concavity of $\log\det$]\label{linear:thm:logdet}
For $A,B\succ0$ and $0\le t\le1$,
\[
\log\det\bigl(tA+(1-t)B\bigr)\ge t\log\det A+(1-t)\log\det B.
\]
\end{theorem}
\begin{proof}
Replacing $A,B$ by $X^*AX,X^*BX$ adds $\log|\det X|^2$ to both sides, so as before we may assume $B=I$ and $A=\operatorname{diag}(d_i)$
with $d_i>0$. The inequality becomes $\sum_i\log(td_i+1-t)\ge\sum_it\log d_i$, which holds term by term because $\log$ is concave:
$\log(td_i+(1-t)\cdot1)\ge t\log d_i+(1-t)\log1$.
\end{proof}

\begin{remark}
For a Gaussian vector with covariance matrix $\Sigma$, the differential entropy is $\frac12\log\det(2\pi e\Sigma)$. Concavity of
$\log\det$ says that mixing covariances increases entropy; Hadamard's inequality says that, for given variances, entropy is largest
when the coordinates are uncorrelated; Fischer's says the same for blocks of coordinates. Information theory turns this around
and proves determinant inequalities from entropy (Cover and Thomas, Chapter~17).
\end{remark}

\section{The Schur product theorem}

The ordinary product of two positive semidefinite matrices need not even be Hermitian. The right product is the entrywise one.

\begin{definition}\label{linear:def:hadamardproduct}
The \emph{Hadamard} (or \emph{Schur}) \emph{product} of two $m\times n$ matrices is $(A\circ B)_{ij}=a_{ij}b_{ij}$.
\end{definition}

\begin{theorem}[Schur, 1911]\label{linear:thm:schurproduct}
If $A,B\succeq0$, then $A\circ B\succeq0$. If $A,B\succ0$, then $A\circ B\succ0$.
\end{theorem}
\begin{proof}
Write $B=\sum_kb_kb_k^*$ (spectral decomposition, $b_k=\sqrt{\lambda_k}v_k$). For one term, $(A\circ b b^*)_{ij}=b_ia_{ij}\overline{b_j}$, that
is, $A\circ bb^*=D_bAD_b^*$ with $D_b=\operatorname{diag}(b)$, which is positive semidefinite by Proposition~\ref{linear:prop:loewner}(i).
Summing over $k$, $A\circ B\succeq0$. If moreover $A,B\succ0$ and $x\ne0$, then $x^*(A\circ B)x=\sum_ky_k^*Ay_k$ with $y_k=D_{b_k}^*x$, whose
$i$-th entry is $\overline{b_{k,i}}x_i$. Pick $i$ with $x_i\ne0$. Since $\sum_k|b_{k,i}|^2=b_{ii}>0$, some $b_{k,i}$ is non-zero, so that
$y_k\ne0$ and $y_k^*Ay_k>0$.
\end{proof}

The Schur product theorem has a striking consequence: entrywise functions with non-negative Taylor coefficients preserve positivity.

\begin{corollary}\label{linear:cor:entrywiseexp}
If $A\succeq0$, then $\bigl(e^{a_{ij}}\bigr)_{i,j}\succeq0$; more generally $\bigl(f(a_{ij})\bigr)\succeq0$ for every power series $f(z)=\sum c_kz^k$
with $c_k\ge0$ converging on the entries. Consequently, for any points $x_1,\dots,x_N\in\R^d$ and $\gamma>0$, the \emph{Gaussian kernel
matrix} $\bigl(e^{-\gamma\|x_i-x_j\|^2}\bigr)_{i,j}$ is positive semidefinite.
\end{corollary}
\begin{proof}
The Schur powers $A^{\circ k}$ are $\succeq0$ by induction ($A^{\circ0}=J=\mathbf 1\mathbf 1^*\succeq0$), so the partial sums
$\sum_{k\le K}c_kA^{\circ k}$ are $\succeq0$, and so is their limit (the set of positive semidefinite matrices is closed). For the kernel,
$e^{-\gamma\|x_i-x_j\|^2}=e^{-\gamma\|x_i\|^2}\,e^{2\gamma\langle x_i,x_j\rangle}\,e^{-\gamma\|x_j\|^2}$: the middle factor is the entrywise
exponential of the Gram matrix $(2\gamma\langle x_i,x_j\rangle)\succeq0$, and the outer factors form a diagonal congruence.
\end{proof}

This is the mathematical engine of \emph{kernel methods} in machine learning: a kernel $k(x,y)$ whose matrices $(k(x_i,x_j))$ are
always positive semidefinite is an inner product in a hidden, possibly infinite-dimensional, feature space, and the Schur product
theorem says that sums, products and limits of such kernels are again kernels. (The Gaussian kernel matrix is in fact positive
definite for distinct points; the proof uses the Fourier transform, Volume~7.)

\section{Gershgorin's discs}

The last tool is the cheapest of all, and it applies to every square matrix, Hermitian or not.

\begin{theorem}[Gershgorin, 1931]\label{linear:thm:gershgorin}
Every eigenvalue of $A\in M_n(\C)$ lies in the union of the discs
\[
D_i=\Bigl\{z\in\C:|z-a_{ii}|\le r_i\Bigr\},\qquad r_i=\sum_{j\ne i}|a_{ij}|.
\]
\end{theorem}
\begin{proof}
Let $Ax=\lambda x$ with $x\ne0$, and pick $i$ with $|x_i|=\max_j|x_j|>0$. The $i$-th equation gives $(\lambda-a_{ii})x_i=\sum_{j\ne
i}a_{ij}x_j$, so $|\lambda-a_{ii}||x_i|\le\sum_{j\ne i}|a_{ij}||x_j|\le r_i|x_i|$.
\end{proof}

\begin{figure}[ht]
\centering
\fig{linear-ch11-gershgorin}
\caption{The four Gershgorin discs of a $4\times4$ complex matrix (centres $a_{ii}$, radii the row sums $r_i$) and its four eigenvalues
(black dots). The isolated disc on the right contains exactly one eigenvalue; the three overlapping discs share the other three.}
\label{linear:fig:gershgorin}
\end{figure}

\begin{corollary}\label{linear:cor:diagdom}
A strictly diagonally dominant matrix ($|a_{ii}|>r_i$ for all $i$) is invertible. A Hermitian matrix with $a_{ii}>r_i$ for all $i$ is
positive definite.
\end{corollary}
\begin{proof}
No disc contains $0$, so $0$ is not an eigenvalue. In the Hermitian case the eigenvalues are real and lie in discs contained in
$\{\operatorname{Re}z>0\}$.
\end{proof}

Applying the theorem to $A^{\mathsf T}$ gives discs with column sums; applying it to $D^{-1}AD$ with a well-chosen positive diagonal $D$
shrinks some discs at the expense of others, a cheap and surprisingly effective way to localise eigenvalues (Exercise~\ref{linear:ex:gershscale}).
A continuity argument, which we do not give, shows more: a union of $k$ discs disjoint from the others contains exactly $k$ eigenvalues,
as in Figure~\ref{linear:fig:gershgorin}.

\section{A gem: Hilbert's inequality and the number $\pi$}

The Hilbert matrix $H=\bigl(\frac1{i+j+1}\bigr)_{i,j\ge0}$ is the Gram matrix of the monomials in $L^2[0,1]$, so its finite sections are
positive definite (Theorem~\ref{linear:thm:gram}), and violently ill-conditioned (Corollary~\ref{linear:cor:hilbert}). Its \emph{largest}
eigenvalues, on the other hand, stay bounded, and the bound is a surprise.

\begin{theorem}[Hilbert's inequality]\label{linear:thm:hilbertineq}
For all real sequences $(a_m)_{m\ge0}$, $(b_n)_{n\ge0}$ with $\sum a_m^2<\infty$ and $\sum b_n^2<\infty$,
\[
\Bigl|\sum_{m,n\ge0}\frac{a_mb_n}{m+n+1}\Bigr|\le\pi\Bigl(\sum_ma_m^2\Bigr)^{1/2}\Bigl(\sum_nb_n^2\Bigr)^{1/2}.
\]
Equivalently, every finite section $H_N$ of the Hilbert matrix has $\|H_N\|_2=\lambda_{\max}(H_N)\le\pi$.
\end{theorem}
\begin{proof}
The idea, due to Issai Schur, is to apply Cauchy--Schwarz with well-chosen weights. Let $p_n=(n+\frac12)^{-1/2}$ and $k_{mn}=\frac1{m+n+1}$.

\emph{Step 1: $\sum_nk_{mn}p_n\le\pi p_m$ for every $m$.} The function $f(x)=\frac1{(m+1+x)\sqrt{x+1/2}}$ is a product of two positive,
decreasing, convex functions on $(-\frac12,\infty)$, hence convex ($(gh)''=g''h+2g'h'+gh''\ge0$). For a convex function, the value at the
midpoint of an interval is at most the average over the interval: $f(n)\le\int_{n-1/2}^{n+1/2}f(x)\,dx$. Summing over $n\ge0$ and
substituting $x+\frac12=y$, then $y=(m+\frac12)u^2$,
\[
\sum_{n\ge0}\frac{p_n}{m+n+1}\le\int_0^\infty\frac{dy}{(m+\frac12+y)\sqrt y}=\frac1{\sqrt{m+\frac12}}\int_0^\infty\frac{2\,du}{1+u^2}=
\pi p_m.
\]

\emph{Step 2: Cauchy--Schwarz.} Write $k_{mn}|a_m||b_n|=\bigl(k_{mn}\frac{p_n}{p_m}\bigr)^{1/2}|a_m|\cdot\bigl(k_{mn}\frac{p_m}{p_n}\bigr)^{1/2}
|b_n|$. Then
\[
\sum_{m,n}k_{mn}|a_m||b_n|\le\Bigl(\sum_{m,n}k_{mn}\frac{p_n}{p_m}a_m^2\Bigr)^{1/2}\Bigl(\sum_{m,n}k_{mn}\frac{p_m}{p_n}b_n^2\Bigr)^{1/2}.
\]
By Step~1, $\sum_nk_{mn}p_n/p_m\le\pi$, so the first factor is at most $(\pi\sum a_m^2)^{1/2}$; by symmetry of $k$, the second is at most
$(\pi\sum b_n^2)^{1/2}$. For the finite sections, $\|H_N\|_2=\max x^*H_Ny$ over unit vectors, which is at most $\pi$ by the inequality.
\end{proof}

The constant $\pi$ cannot be improved, as Schur showed in 1911: $\lambda_{\max}(H_N)\to\pi$, but with a gap that de Bruijn and Wilf
showed in 1962 to behave like $\frac{\pi^5}{2\log^2N}$: so slowly that at $N=3000$ the largest eigenvalue is still only about $2.53$
(a numerical calculation watches it crawl). Hilbert proved his inequality with the
constant $2\pi$ in his lectures on integral equations, published by Hermann Weyl in 1908; the integral version,
$\iint\frac{f(x)g(y)}{x+y}\,dx\,dy\le\pi\|f\|_2\|g\|_2$, is a cornerstone of Hardy, Littlewood and P\'olya's book \emph{Inequalities}.

\begin{beauty}
Numbers are compared by size; matrices are compared by ellipsoids. Under that one change of viewpoint, volumes obey Hadamard,
blocks obey Fischer, sums obey Minkowski and mixtures obey the concavity of $\log\det$, each proved by reducing a pair of matrices to a
pair of diagonal ones and quoting the inequality between arithmetic and geometric means. And when the matrix is Hilbert's, the
best constant in the comparison is $\pi$, pulled out of an integral by a convex function.
\end{beauty}

\section{Exercises}

\begin{exercise}[$\star$]\label{linear:ex:commutingsquare}
Show that if $A\succeq B\succeq0$ and $AB=BA$, then $A^2\succeq B^2$.
\end{exercise}

\begin{exercise}[$\star$]\label{linear:ex:blockcs}
Show that if $\bigl(\begin{smallmatrix}A&B\\B^*&C\end{smallmatrix}\bigr)\succeq0$, then $|x^*By|^2\le(x^*Ax)(y^*Cy)$ for all vectors $x,y$.
\end{exercise}

\begin{exercise}[$\star\star$]\label{linear:ex:sherman}
Let $A$ be invertible and $1+v^*A^{-1}u\ne0$. Show that $\det(A+uv^*)=\det A\,(1+v^*A^{-1}u)$ and that
\[
(A+uv^*)^{-1}=A^{-1}-\frac{A^{-1}uv^*A^{-1}}{1+v^*A^{-1}u}
\]
(Sherman--Morrison). How many operations does it take to update $A^{-1}$ after a rank-one change?
\end{exercise}

\begin{exercise}[$\star\star$]\label{linear:ex:hadamard4}
Let $H$ be an $n\times n$ Hadamard matrix with $n\ge3$. Normalise its first three rows to begin with $1$s by multiplying columns by
$-1$, and count the columns of each sign pattern in the first three rows. Deduce that $4\mid n$.
\end{exercise}

\begin{exercise}[$\star\star$]\label{linear:ex:maxdet3}
Show that the determinant of an $n\times n$ matrix of $\pm1$ is divisible by $2^{n-1}$. Deduce that the maximum determinant of a $3\times3$
matrix of $\pm1$ is $4$, although Hadamard's bound is $3^{3/2}\approx5.2$, and find a matrix that attains it.
\end{exercise}

\begin{exercise}[$\star$]\label{linear:ex:minkowskieq}
Show that equality holds in Minkowski's inequality for $A,B\succ0$ if and only if $B=cA$ for some $c>0$.
\end{exercise}

\begin{exercise}[$\star\star$]\label{linear:ex:logdetderiv}
For $A\succ0$ and Hermitian $H$, show that $\frac d{dt}\log\det(A+tH)\big|_{t=0}=\operatorname{tr}(A^{-1}H)$ and
$\frac{d^2}{dt^2}\log\det(A+tH)\big|_{t=0}=-\operatorname{tr}\bigl((A^{-1}H)^2\bigr)\le0$. This is a second proof of Theorem~\ref{linear:thm:logdet}.
\end{exercise}

\begin{exercise}[$\star\star$]\label{linear:ex:schurmin}
Show that for $A,B\succeq0$, $\lambda_{\min}(A\circ B)\ge\lambda_{\min}(A)\min_ib_{ii}$.
\end{exercise}

\begin{exercise}[$\star\star$]\label{linear:ex:kernels}
Show that the following matrices are positive semidefinite, for any choice of the points: $\bigl(\min(s_i,s_j)\bigr)$ for $s_i\ge0$;
$\bigl(\frac1{x_i+x_j}\bigr)$ for $x_i>0$; $\bigl(\cos(\theta_i-\theta_j)\bigr)$ for real $\theta_i$; and $\bigl(\frac1{1-\langle u_i,u_j\rangle}\bigr)$
for vectors $u_i\in\R^d$ with $\|u_i\|<1$.
\end{exercise}

\begin{exercise}[$\star$]\label{linear:ex:gershscale}
Let $A=\begin{pmatrix}1&\varepsilon&\varepsilon\\\varepsilon&2&\varepsilon\\\varepsilon&\varepsilon&3\end{pmatrix}$ with $0<\varepsilon<\frac14$. Weyl's
inequality (Corollary~\ref{linear:cor:weylpert}) puts the smallest eigenvalue $\mu$ of $A$ within $2\varepsilon$ of $1$. Use Gershgorin's theorem
for $D^{-1}AD$, with a well-chosen diagonal $D$, to show that in fact $|\mu-1|\le8\varepsilon^2$.
\end{exercise}

\begin{exercise}[$\star\star$]\label{linear:ex:diagdomdet}
Show that a strictly diagonally dominant real matrix with positive diagonal has a positive determinant. (Deform it to its diagonal.)
\end{exercise}

\begin{exercise}[$\star\star\star$]\label{linear:ex:oppenheim}
(Oppenheim's inequality.) For $A,B\succeq0$, show that $\det(A\circ B)\ge\det A\prod_ib_{ii}$. (Induct on $n$, using Schur complements
and the Schur product theorem.)
\end{exercise}

\section{Hints}

\begin{hint}{linear:ex:commutingsquare}
Commuting Hermitian matrices are diagonal in a common orthonormal basis (Exercise~\ref{linear:ex:simultaneous}).
\end{hint}
\begin{hint}{linear:ex:blockcs}
Look at the $2\times2$ matrix obtained by testing $M$ against $(sx,ty)$, or apply Proposition~\ref{linear:prop:schurcomplement} to it.
\end{hint}
\begin{hint}{linear:ex:sherman}
Factor $A+uv^*=A(I+A^{-1}uv^*)$ and use Proposition~\ref{linear:prop:detIAB}. For the inverse, multiply out.
\end{hint}
\begin{hint}{linear:ex:hadamard4}
Let $a,b,c,d$ be the numbers of columns whose entries in rows $2,3$ are $(+,+),(+,-),(-,+),(-,-)$. Orthogonality of rows $1,2$, $1,3$ and
$2,3$ gives three linear equations.
\end{hint}
\begin{hint}{linear:ex:maxdet3}
Add the first row to the others: the new rows have even entries.
\end{hint}
\begin{hint}{linear:ex:minkowskieq}
Look at the equality case in AM--GM in the proof.
\end{hint}
\begin{hint}{linear:ex:logdetderiv}
$\log\det(A+tH)=\log\det A+\sum_i\log(1+t\mu_i)$, where $\mu_i$ are the eigenvalues of $A^{-1/2}HA^{-1/2}$.
\end{hint}
\begin{hint}{linear:ex:schurmin}
$A-\lambda_{\min}(A)I\succeq0$.
\end{hint}
\begin{hint}{linear:ex:kernels}
Find Gram matrices: $\min(s,t)=\int_0^\infty\mathbf 1_{[0,s]}\mathbf 1_{[0,t]}$, $\frac1{x+y}=\int_0^\infty e^{-xr}e^{-yr}\,dr$; for the last one,
a geometric series.
\end{hint}
\begin{hint}{linear:ex:gershscale}
Try $D=\operatorname{diag}(1,\delta,\delta)$ and choose $\delta$ so that the first disc is small while the other two stay away from $\mu$.
\end{hint}
\begin{hint}{linear:ex:diagdomdet}
$A_t=D+t(A-D)$ stays strictly diagonally dominant for $t\in[0,1]$, so $\det A_t$ never vanishes.
\end{hint}
\begin{hint}{linear:ex:oppenheim}
Write $A=\bigl(\begin{smallmatrix}a&\alpha^*\\\alpha&A_1\end{smallmatrix}\bigr)$ and note that $A_1-\alpha\alpha^*/a\succeq0$.
\end{hint}

\section{Solutions}

\begin{solution}{linear:ex:commutingsquare}
\idea Commuting matrices can be diagonalised together, and then the question is about numbers.
\writeup By Exercise~\ref{linear:ex:simultaneous}, $A=U\operatorname{diag}(\alpha_i)U^*$ and $B=U\operatorname{diag}(\beta_i)U^*$ with the same
unitary $U$. Then $A\succeq B\succeq0$ means $\alpha_i\ge\beta_i\ge0$, so $\alpha_i^2\ge\beta_i^2$, which is $A^2\succeq B^2$. The
counterexample of the text shows that commutativity cannot be dropped.
\end{solution}

\begin{solution}{linear:ex:blockcs}
\idea A positive semidefinite matrix stays positive semidefinite when compressed to a $2\times2$ matrix, and a $2\times2$ positive
semidefinite matrix has non-negative determinant.
\writeup Let $Y$ be the matrix with the two columns $(x,0)$ and $(0,y)$. By Proposition~\ref{linear:prop:loewner}(i), $Y^*MY=\bigl(\begin{smallmatrix}
x^*Ax&x^*By\\y^*B^*x&y^*Cy\end{smallmatrix}\bigr)\succeq0$, so its determinant $(x^*Ax)(y^*Cy)-|x^*By|^2$ is $\ge0$.
\end{solution}

\begin{solution}{linear:ex:sherman}
\idea Factor out $A$ and use $\det(I+ab^*)=1+b^*a$; for the inverse, guess a rank-one correction and verify.
\writeup $\det(A+uv^*)=\det A\det(I+(A^{-1}u)v^*)=\det A\,(1+v^*A^{-1}u)$ by Proposition~\ref{linear:prop:detIAB}. With $c=1+v^*A^{-1}u$,
\[
\begin{aligned}
(A+uv^*)\Bigl(A^{-1}-\frac{A^{-1}uv^*A^{-1}}c\Bigr)&=I+uv^*A^{-1}-\frac{uv^*A^{-1}+u(v^*A^{-1}u)v^*A^{-1}}c\\
&=I+uv^*A^{-1}-\frac{1+v^*A^{-1}u}c\,uv^*A^{-1}=I.
\end{aligned}
\]
Knowing $A^{-1}$, the update costs two matrix-vector products and one outer product: $O(n^2)$ operations instead of $O(n^3)$.
\end{solution}

\begin{solution}{linear:ex:hadamard4}
\idea Orthogonality of $\pm1$ vectors is a counting statement: two rows agree in exactly half of the positions.
\writeup Multiplying columns by $-1$ preserves $HH^{\mathsf T}=nI$, so assume the first row is all $1$s. Rows $2$ and $3$ are orthogonal to
row $1$ and to each other. With $a,b,c,d$ as in the hint, $a+b+c+d=n$; orthogonality with row $1$ gives $a+b=c+d$ (row $2$) and $a+c=
b+d$ (row $3$); orthogonality of rows $2,3$ gives $a+d=b+c$. Subtracting the first two equations gives $b-c=c-b$, so $b=c$; adding them gives $a=d$; then
the third reads $2a=2b$. So $a=b=c=d$ and $n=4a$.
\end{solution}

\begin{solution}{linear:ex:maxdet3}
\idea Row operations do not change the determinant and can make every entry outside the first row even.
\writeup Adding the first row to every other row gives rows with entries in $\{0,\pm2\}$; factoring $2$ out of each of the $n-1$ rows
leaves an integer matrix. So $2^{n-1}\mid\det$. For $n=3$, $|\det|$ is a multiple of $4$ and at most $3^{3/2}<8$, so at most $4$. The matrix
\[
\begin{pmatrix}1&1&1\\1&-1&1\\1&1&-1\end{pmatrix}
\]
has determinant $1(1-1)-1(-1-1)+1(1+1)=4$. Hadamard's bound is not attained, as it must not be: $3$ is not a multiple of $4$.
\end{solution}

\begin{solution}{linear:ex:minkowskieq}
\idea Every step of the proof was AM--GM, and AM--GM is an equality only when all terms are equal.
\writeup After the reduction $A=I$, $B=\operatorname{diag}(d_i)$ with $d_i>0$, equality requires equality in both AM--GM inequalities,
so $\frac{1}{1+d_i}$ is the same for all $i$: all $d_i$ equal some $c>0$, and $B=cI=cA$. Undoing the congruence, $B=cA$. Conversely,
$\det((1+c)A)^{1/n}=(1+c)\det(A)^{1/n}$.
\end{solution}

\begin{solution}{linear:ex:logdetderiv}
\idea Reduce to eigenvalues, where $\log\det$ is a sum of logarithms.
\writeup $A+tH=A^{1/2}(I+tK)A^{1/2}$ with $K=A^{-1/2}HA^{-1/2}$ Hermitian, eigenvalues $\mu_i$. So $\log\det(A+tH)=\log\det A+\sum_i\log(1+
t\mu_i)$, with first derivative $\sum\mu_i=\operatorname{tr}K=\operatorname{tr}(A^{-1}H)$ and second derivative $-\sum\mu_i^2=-\operatorname{tr}
K^2=-\operatorname{tr}((A^{-1}H)^2)\le0$ at $t=0$. A function on a convex set whose second derivative along every segment is $\le0$ is
concave; this applies at every point $A$ of the open cone, and gives Theorem~\ref{linear:thm:logdet} again.
\end{solution}

\begin{solution}{linear:ex:schurmin}
\idea Peel off the smallest eigenvalue of $A$ as a multiple of the identity, which the Schur product turns into a diagonal matrix.
\writeup Let $\lambda=\lambda_{\min}(A)\ge0$. Then $A-\lambda I\succeq0$, so $(A-\lambda I)\circ B\succeq0$ by Theorem~\ref{linear:thm:schurproduct},
that is, $A\circ B\succeq\lambda I\circ B=\lambda\operatorname{diag}(b_{ii})\succeq\lambda\min_ib_{ii}\,I$. By Proposition~\ref{linear:prop:loewner}(ii),
$\lambda_{\min}(A\circ B)\ge\lambda\min_ib_{ii}$.
\end{solution}

\begin{solution}{linear:ex:kernels}
\idea A matrix of inner products is positive semidefinite, whatever the inner product space; so it suffices to write each entry as
an inner product, possibly in a function space.
\writeup (a) With $f_i=\mathbf 1_{[0,s_i]}$,
\[
\min(s_i,s_j)=\int_0^\infty f_i(r)f_j(r)\,dr=\langle f_i,f_j\rangle_{L^2},
\]
so $\sum_{i,j}c_ic_j
\min(s_i,s_j)=\|\sum c_if_i\|^2\ge0$. (This is the covariance of Brownian motion.) (b) $\frac1{x_i+x_j}=\int_0^\infty e^{-x_ir}e^{-x_jr}\,dr$, the
same argument. (c) $\cos(\theta_i-\theta_j)=\cos\theta_i\cos\theta_j+\sin\theta_i\sin\theta_j$, the Gram matrix of the vectors
$(\cos\theta_i,\sin\theta_i)$. (d) $\frac1{1-\langle u_i,u_j\rangle}=\sum_{k\ge0}\langle u_i,u_j\rangle^k$, a convergent series with non-negative
coefficients applied entrywise to a Gram matrix (it converges since $|\langle u_i,u_j\rangle|<1$ by Cauchy--Schwarz):
Corollary~\ref{linear:cor:entrywiseexp}.
\end{solution}

\begin{solution}{linear:ex:gershscale}
\idea A diagonal similarity does not change the eigenvalues but moves off-diagonal weight between rows and columns. Shrinking
the off-diagonal part of row $1$ shrinks its disc; the price is paid by the other rows, whose discs only need to stay away from $\mu$.
\writeup First, $A=\operatorname{diag}(1,2,3)+E$ where $E$ has zero diagonal and off-diagonal entries $\varepsilon$; its eigenvalues are $2\varepsilon$
and $-\varepsilon$ (twice), so $\|E\|_2=2\varepsilon$, and by Weyl the eigenvalues of $A$ are within $2\varepsilon<\frac12$ of $1,2,3$. In
particular $\mu\le1+2\varepsilon<\frac32$. Now let $D=\operatorname{diag}(1,\delta,\delta)$ with $\delta=4\varepsilon$. The entries of
$D^{-1}AD$ are $a_{ij}d_j/d_i$, so its rows are
\[
(1,\ \varepsilon\delta,\ \varepsilon\delta),\qquad(\varepsilon/\delta,\ 2,\ \varepsilon),\qquad(\varepsilon/\delta,\ \varepsilon,\ 3).
\]
The discs have centres $1,2,3$ and radii $2\varepsilon\delta=8\varepsilon^2$, $\frac\varepsilon\delta+\varepsilon=\frac14+\varepsilon$ and
$\frac14+\varepsilon$. The second and third discs lie in $\{\operatorname{Re}z\ge2-\frac14-\varepsilon\}\subseteq\{\operatorname{Re}z>\frac32\}$,
so they do not contain $\mu$. By Gershgorin, $\mu$ lies in the first disc: $|\mu-1|\le8\varepsilon^2$. (Second-order perturbation theory
predicts $\mu\approx1-\varepsilon^2-\frac{\varepsilon^2}2$, so the true distance is about $1.5\,\varepsilon^2$; a direct computation verifies it.)
\end{solution}

\begin{solution}{linear:ex:diagdomdet}
\idea The determinant is a continuous function that never vanishes along a path from the matrix to its diagonal, so it keeps its sign.
\writeup Let $D$ be the diagonal part of $A$ and $A_t=D+t(A-D)$, $t\in[0,1]$. The diagonal of $A_t$ is that of $A$, and its off-diagonal
row sums are $tr_i\le r_i<a_{ii}$, so $A_t$ is strictly diagonally dominant and invertible (Corollary~\ref{linear:cor:diagdom}): $\det A_t\ne0$. The
function $t\mapsto\det A_t$ is continuous (a polynomial) and never zero on $[0,1]$, so it has constant sign, that of $\det D=\prod a_{ii}>0$.
\end{solution}

\begin{solution}{linear:ex:oppenheim}
\idea Peel off the first row and column with Schur complements on both sides; the Schur product theorem shows that the complement
for $A\circ B$ dominates the Schur product of the complement for $A$ with a block of $B$, and induction does the rest.
\writeup By continuity (replace $A$ by $A+\varepsilon I$ and let $\varepsilon\to0$) we may assume $A\succ0$; the case $n=1$ is an equality.
Write
\[
A=\begin{pmatrix}a&\alpha^*\\\alpha&A_1\end{pmatrix},\qquad B=\begin{pmatrix}b&\beta^*\\\beta&B_1\end{pmatrix},
\]
so $a>0$, and $S=A_1-\alpha\alpha^*/a\succ0$ with $\det A=a\det S$ (Proposition~\ref{linear:prop:schurcomplement}). If $b=0$, then $\beta=0$ (a
positive semidefinite matrix with a zero diagonal entry has a zero row), so $A\circ B$ has a zero row and both sides vanish. If $b>0$,
then $B_1-\beta\beta^*/b\succeq0$, and the Schur complement of the entry $ab$ in $A\circ B$ is
\[
A_1\circ B_1-\frac{(\alpha\circ\beta)(\alpha\circ\beta)^*}{ab}=S\circ B_1+\frac1a\,\alpha\alpha^*\circ\Bigl(B_1-\frac{\beta\beta^*}b\Bigr),
\]
as one sees by expanding the right side and using $(\alpha\alpha^*)\circ(\beta\beta^*)=(\alpha\circ\beta)(\alpha\circ\beta)^*$. The last term is
$\succeq0$ by the Schur product theorem, so this complement is $\succeq S\circ B_1\succeq0$. By Proposition~\ref{linear:prop:loewner}(ii) and
induction (applied to $S$ and $B_1$),
\[
\begin{aligned}
\det(A\circ B)&=ab\,\det\Bigl(A_1\circ B_1-\frac{(\alpha\circ\beta)(\alpha\circ\beta)^*}{ab}\Bigr)\\
&\ge ab\,\det(S\circ B_1)\ge ab\,\det S\prod_{i\ge2}b_{ii}=\det A\prod_ib_{ii}.
\end{aligned}
\]
\end{solution}

\section*{Sources and further reading}

\begin{itemize}
\item J.~Hadamard, \emph{R\'esolution d'une question relative aux d\'eterminants}, Bulletin des Sciences Math\'ematiques 17 (1893),
240--246; H.~Kharaghani and B.~Tayfeh-Rezaie, \emph{A Hadamard matrix of order 428}, Journal of Combinatorial Designs 13 (2005), 435--440.
\item R.~Bhatia, \emph{Matrix Analysis}, Springer, 1997, Chapters V (operator monotone functions, L\"owner's theorem) and IX; R.~Bhatia,
\emph{Positive Definite Matrices}, Princeton University Press, 2007.
\item R.~A.~Horn and C.~R.~Johnson, \emph{Matrix Analysis}, 2nd ed., Cambridge University Press, 2013, Chapters 6 (Gershgorin) and 7
(positive definite matrices, Hadamard, Fischer, Oppenheim, Schur product).
\item I.~Schur, \emph{Bemerkungen zur Theorie der beschr\"ankten Bilinearformen mit unendlich vielen Ver\"anderlichen}, Journal f\"ur die
reine und angewandte Mathematik 140 (1911), 1--28 (Schur product theorem; Hilbert's inequality with the sharp constant $\pi$).
\item G.~H.~Hardy, J.~E.~Littlewood and G.~P\'olya, \emph{Inequalities}, Cambridge University Press, 1934, Chapter IX (Hilbert's inequality);
J.~M.~Steele, \emph{The Cauchy--Schwarz Master Class}, Cambridge University Press, 2004, Chapter 10.
\item S.~Gershgorin, \emph{\"Uber die Abgrenzung der Eigenwerte einer Matrix}, Izvestiya Akademii Nauk SSSR 6 (1931), 749--754; R.~S.~Varga,
\emph{Ger\v{s}gorin and His Circles}, Springer, 2004.
\item K.~L\"owner, \emph{\"Uber monotone Matrixfunktionen}, Mathematische Zeitschrift 38 (1934), 177--216; E.~Heinz, \emph{Beitr\"age zur
St\"orungstheorie der Spektralzerlegung}, Mathematische Annalen 123 (1951), 415--438.
\item B.~Sch\"olkopf and A.~J.~Smola, \emph{Learning with Kernels}, MIT Press, 2002, Chapter 2 (positive definite kernels).
\item T.~M.~Cover and J.~A.~Thomas, \emph{Elements of Information Theory}, 2nd ed., Wiley, 2006, Chapter 17 (determinant inequalities
from entropy).
\end{itemize}
